\subsection{負AS2平面を別々に正規化する正量子二元}
\label{subsec:h3-order2-22-rank4-opposite-separated}
\begin{theorem}\label{thm:h3-order2-22-rank4-opposite-separated}
負AS2平面$T=k\langle u,v\rangle/(r)$を、正生成元$x,y$がそれぞれ自己同型$F,G$で正規化し、
正関係が$xy=pyx$、$p\ne0$となるgraph候補を考える。
中山半単純部分が$(+,+,-,-)$なら、全Jordan型を含めたこのchartのH3部分も
定理\ref{thm:fixed-sign-classification}の固定符号包絡の切断だけである。
正平面がJordan型のopposite-separated表示は、この定理に含めない。
\end{theorem}
\begin{proof}
PBWの結合則から$FG=GF$であり、二回のOre公式は
\[
 \mu_B|_T=(FG)^{-1}\mu_T,\qquad
 \mu_B(x)=p\operatorname{hdet}(F)x,\qquad
 \mu_B(y)=p^{-1}\operatorname{hdet}(G)y
\]
となる。正固有値の式は厳密に1である。
必要な三次のPBW表示は
\[
 G_3=H_3(x,y)+Q_x(u,v)x+Q_y(u,v)y
\]
である。$Q_x,Q_y\in T_2$とし、負平面上の正常化作用を
$\phi=\mu_T^{-1}FG$と書く。
その半単純部分は$-I$である。

$T$が可換なら、$Q_x\ne0$は$Q_xF(t)=\phi(t)Q_x$から$G=I$を、
$Q_y\ne0$は$F=I$を強制する。
 $F$も$G$も恒等でなければ、$G_3$は正量子平面の斉次三次式だけとなり、一次因子へ分解する。
$G=I$の場合、$F$の半単純部分は$-I$、$\det F=1$である。
正中山の式から$p=1$となる。
$H_3$中の$x^iy^j$が負平面を正常化するには、$F^i=\phi=F$が必要であり、
半単純部分を比較して$i$は奇数となる。
従って$H_3$は$x$で割れ、残る項も$Q_xx$だけなので、全$G_3$が$x$で割れる。
$F=I$の場合も同様に$y$が共通因子となる。

$T$が量子平面$uv=qvu$、$q\ne1$の場合を考える。
$q\ne-1$では$F,G$は対角なので、包絡全体がskew多項式環となり、既存の全skew分類を適用できる。
$q=-1$でも両方対角なら同じである。
両方反対角の場合、$\mu_T=-I$と中山の半単純部分の条件から$FG=I$となる。
片方だけが反対角だと負平面の二固有値が相反するので仮定に反する。

両反対角の場合、$H_3$の各単項式の正常化作用は$F^{i-j}$で、$i+j=3$である。
奇数乗の反対角作用は$\phi=-I$に一致しないので、$H_3=0$である。
$Q_xF(t)=-tQ_x$が非零なら、平方を取って
\[
 Q_xF(u^2)=u^2Q_x
\]
を得る。$u^2,v^2$は中心であるから、整域性から$F(u^2)=u^2$が必要である。
ところが反対角$F$は$u^2$を$v^2$の非零倍へ移すので矛盾する。
$Q_y$も同様に零であり、この場合は正常核が零となる。
従って量子負平面の場合の全H3は、skew分類にある固定符号包絡だけである。

最後にJordan負平面$vu-uv=\kappa u^2$、$\kappa\ne0$を調べる。
$F(u)=a u,G(u)=c u$と書けば、負中山の条件から$ac=-1$である。
$Q_xF(u)=\phi(u)Q_x=-uQ_x$をPBWの最高$v$次数で比較すると、
非零$Q_x$には$a=-1$が必要となる。
同様に非零$Q_y$には$c=-1$が必要である。
どちらも成り立たなければ、$G_3$は正量子平面の三次式だけであり、積に分解する。

$a=-1$なら$c=1$で、$\operatorname{hdet}F=a^2$と正中山から$p=1$を得る。
$H_3$の各単項式の作用$F^iG^j$と$\phi$の半単純部分を比較すると、$i$は奇数である。
従って全$H_3$が$x$で割れ、$Q_y=0$なので全$G_3$も$x$で割れる。
$c=-1$の場合は同様に$y$で割れる。
全ての因数分解は一次と二次の非零因子を与え、H3の整域商に反する。
これで負AS2平面の三標準形を尽くし、主張を得る。
\end{proof}

検証器は
\[
 \texttt{verify\_cubic\_order2\_22\_rank4\_opposite\_separated.py}
\]
である。Jordan二次係数の絡み合わせ行列の行列式$(a+1)^3$、
反対角作用の中心平方への作用、奇数乗と例外係数を確認する。
このchart認識を、全mixed graphの被覆や族の降下と取り違えていない。
